\section{Observability computations for SIR, Lotka--Volterra, and OGTT}
\label{app:observability}

This appendix documents the calculations and code used to support the observability/identifiability statements in
Section~\ref{sec:identifiability}.
We treat parameters as augmented constant states and follow the nonlinear observability framework of Hermann and Krener~\cite{hermann1977nonlinear}.

\subsection{SIR model: Lie-derivative-based observability matrix}
\label{app:obs-sir}

We take the augmented state
\[
z = (S,I,R,\beta,\gamma),
\]
with observation $h(z)=I$.
The vector field for the augmented system is
\[
\Sigma(z) = (\dot{S},\dot{I},\dot{R},0,0)
=
(-\beta S I,\ \beta S I-\gamma I,\ \gamma I,\ 0,\ 0).
\]

We compute Lie derivatives of $h$ along $\Sigma$ and their gradients:
\begin{align*}
L_\Sigma^0 h(z) &= h(z) = I,\\
L_\Sigma^1 h(z) &= \nabla h(z)\cdot \Sigma(z) = \beta S I - \gamma I,\\
L_\Sigma^2 h(z) &= \nabla L_\Sigma^1 h(z)\cdot \Sigma(z),\\
&\vdots
\end{align*}
and assemble the observation map
\[
\Phi(z)
=
\begin{bmatrix}
L_\Sigma^0 h(z)\\
L_\Sigma^1 h(z)\\
\vdots\\
L_\Sigma^{4} h(z)
\end{bmatrix},
\qquad
\mathbf{O}(z) := \nabla_z \Phi(z)\in\mathbb{R}^{5\times 5}.
\]

Symbolic computation (e.g., using \texttt{sympy}) yields
\[
\mathrm{rank}\,\mathbf{O}(z)=5
\]
for generic $z$, confirming local observability from infected-only observations.
(Explicit expressions for $\mathbf{O}(z)$ are omitted for brevity but can be generated by the code provided in our repository.)

\subsection{Lotka--Volterra model: symbolic rank deficiency}
\label{app:obs-lv}

For the Lotka--Volterra system we augment the state as
\[
z = (x,y,\alpha,\beta,\gamma,\delta),
\]
with vector field
\[
\Sigma(z) =
(\alpha x - \beta x y,\ \delta x y - \gamma y,\ 0,\ 0,\ 0,\ 0),
\]
and observation $h(z)=x$.

The following \texttt{sympy} script computes the observation matrix and its rank:

\begin{verbatim}
import sympy as sp

x, y, a, b, g, d = sp.symbols('x y a b g d')
z = sp.Matrix([x, y, a, b, g, d])

sigma = sp.Matrix([
    a*x - b*x*y,
    d*x*y - g*y,
    0, 0, 0, 0
])

h = x

Phi = []
current_L = h
Phi.append(current_L)

for _ in range(len(z)):
    grad_L = sp.Matrix([current_L]).jacobian(z)
    current_L = (grad_L * sigma)[0]
    Phi.append(current_L)

Phi_mat = sp.Matrix(Phi)
O = Phi_mat.jacobian(z)  # observation matrix
rank = O.rank()
print("rank(O) =", rank)
\end{verbatim}

Running this code yields \texttt{rank(O) = 5}, while the augmented state dimension is $6$.
Thus the observation matrix is rank-deficient and the Lotka--Volterra model is not locally observable from prey-only observations.
Figure~\ref{fig:lv_unidentifiable} in the main text provides a concrete example of two distinct states/parameters producing essentially identical $x(t)$.

\subsection{OGTT model: Jacobian-based rank estimation}
\label{app:obs-ogtt}

For OGTT, the augmented state and vector field are given in Section~\ref{subsec:ident-ogtt}.
Direct symbolic construction of the observation matrix $\mathbf{O}(z)$ is intractable, so we estimate its rank numerically using the normalized
Jacobian method described in Section~\ref{subsec:ident-ogtt}.

A simplified implementation sketch is:

\begin{verbatim}
def estimate_rank(system, z0, times, eps=1e-4, delta=1e-3):
    G_nom = simulate_glucose(system, z0, times)  # G_nom(t)
    d = len(z0)
    J = np.zeros((len(times), d))

    for k in range(d):
        dz = np.zeros_like(z0)
        dz[k] = z0[k] * eps
        z_pert = z0 + dz
        G_pert = simulate_glucose(system, z_pert, times)

        s_raw = (G_pert - G_nom) / dz[k]      # finite-difference sensitivity
        s_norm = normalize_sensitivity(s_raw) # state-dependent rescaling
        J[:, k] = s_norm

    U, S, Vt = np.linalg.svd(J, full_matrices=False)
    r = np.sum(S > delta)
    return r, S
\end{verbatim}

Here \texttt{simulate\_glucose} integrates the OGTT dynamics and returns the glucose component $G(t)$ at the specified sampling times,
and \texttt{normalize\_sensitivity} rescales sensitivities to account for heterogeneous state scales (see Section~\ref{subsec:ident-ogtt}).
Using this procedure with sufficiently many time points, we obtain six singular values all well above machine precision, leading to the
rank estimate $r=6$ and indicating full-rank observability under idealized conditions.